Put \(h=\min\{B,n-B\}=nq\).  For one row, put \(\xi=Z_{t\cdot}-p\mathbf 1_n\).  If \(B\le n/2\), then
\[
 \xi^{\mathsf T}u=\sum_{j\in R}(u_j-\bar u),
 \qquad \bar u=n^{-1}\sum_{j=1}^n u_j,
 \tag{1}
\]
where \(R\) is a uniform \(h\)-subset.  If \(B>n/2\), the same identity holds with a minus sign and \(R\) the untreated \(h\)-subset.  Thus it is enough to analyze the right side of (1).

We first repair the moment calculation by an exchangeable-pair recursion.  Let \(a\in\mathbb R^n\) satisfy \(\sum_j a_j=0\), let \(W=\sum_{j\in R}a_j\), choose \(I\) uniformly from \(R\) and \(J\) uniformly from \(R^c\), and put
\[
 R'=(R\setminus\{I\})\cup\{J\},
 \qquad W'=W-a_I+a_J,
 \qquad \lambda_0=\frac{n}{h(n-h)}.
 \tag{2}
\]
The pair \((W,W')\) is exchangeable, and direct conditional averaging gives
\[
 \mathbb E(W'-W\mid R)
 =-\left(\frac1h+\frac1{n-h}\right)W
 =-\lambda_0W.
 \tag{3}
\]
For every integer \(\ell\ge1\), exchangeability and (3) imply the exact identity
\[
 \mathbb EW^{2\ell}
 =\frac1{2\lambda_0}\,
 \mathbb E\!\left[(W'-W)
       (W'^{\,2\ell-1}-W^{2\ell-1})\right].
 \tag{4}
\]
Indeed, exchangeability changes the sign of
\((W'-W)\{W'^{\,2\ell-1}+W^{2\ell-1}\}\), so its expectation is zero, while conditioning
\((W'-W)W^{2\ell-1}\) on \(R\) gives
\(-\lambda_0\mathbb EW^{2\ell}\).

The conditional squared increment has the deterministic bound
\[
 \begin{aligned}
 \mathbb E\!\left[(W'-W)^2\mid R\right]
 &\le
 2h^{-1}\sum_{i\in R}a_i^2+
 2(n-h)^{-1}\sum_{j\notin R}a_j^2\\
 &\le2\lambda_0\lVert a\rVert_2^2.
 \end{aligned}
 \tag{5}
\]
Using
\[
 |(x-y)(x^{2\ell-1}-y^{2\ell-1})|
 \le(2\ell-1)(x-y)^2
       (|x|^{2\ell-2}+|y|^{2\ell-2})
\]
in (4), then exchangeability and (5), yields the recursion
\[
 \mathbb E|W|^{2\ell}
 \le2(2\ell-1)\lVert a\rVert_2^2
       \mathbb E|W|^{2\ell-2}.
 \tag{6}
\]
The exact covariance of simple random sampling gives
\[
 \mathbb EW^2
 =\frac{h(n-h)}{n(n-1)}\lVert a\rVert_2^2
 \le2q\lVert a\rVert_2^2.
\]
Iteration of (6), and Cauchy--Schwarz between the two adjacent even moments for odd orders, therefore prove, for every integer \(r\ge2\),
\[
 \mathbb E|W|^r
 \le q(C_1r)^{r/2}\lVert a\rVert_2^r.
 \tag{7}
\]
Thus the relevant \(2\ell\)-th moment is \(q(C\ell)^\ell\), not the invalid
\((2\ell)!\) bound left after exponential-series division.

Since \(a=u-\bar u\mathbf 1_n\) satisfies \(\lVert a\rVert_2\le\lVert u\rVert_2\), (7) implies the Bernstein moment condition.  More explicitly, with \(V=W^2-\mathbb EW^2\), there are numerical \(v_0,c_0'>0\) such that, for every \(\ell\ge2\),
\[
 \mathbb E|V|^\ell
 \le \frac{\ell!}{2}(v_0q\lVert u\rVert_2^4)
       (c_0'\lVert u\rVert_2^2)^{\ell-2}.
 \tag{8}
\]
The same bound (7), applied to both signs of the exponential parameter and using \(\mathbb EW=0\), gives the analogous Bernstein moment-generating-function bound for \(W\), with variance proxy \(Cq\lVert u\rVert_2^2\) and scale \(C\lVert u\rVert_2\).  Independence of the rows under \(\Pi_{m,B}\), the exponential Markov inequality, and (8) therefore yield, for every \(x>0\),
\[
 \Pr\!\left\{
  \left|\frac1m\sum_{t=1}^m(\xi_t^{\mathsf T}u)^2
        -\mathbb E(\xi^{\mathsf T}u)^2\right|
  >C_2\lVert u\rVert_2^2
       \left(\sqrt{\frac{qx}{m}}+\frac{x}{m}\right)
 \right\}\le2e^{-x},
 \tag{9}
\]
and
\[
 \Pr\!\left\{
  |\bar\xi^{\mathsf T}u|>
  C_2\lVert u\rVert_2
       \left(\sqrt{\frac{qx}{m}}+\frac{x}{m}\right)
 \right\}\le2e^{-x},
 \qquad \bar\xi=m^{-1}\sum_{t=1}^m\xi_t.
 \tag{10}
\]
Thus (9), unlike the discarded argument, follows from a genuine Bernstein moment bound for the quadratic variable.

Let
\[
 M=m^{-1}\sum_{t=1}^m\xi_t\xi_t^{\mathsf T},
 \qquad \Sigma=\mathbb E(\xi\xi^{\mathsf T}).
\]
For each support of size at most \(r\), choose a \(1/4\)-net of its Euclidean unit sphere with at most \(9^r\) points.  Padding smaller supports to size \(r\) shows that the number of choices is at most \((en/r)^r\).  Applying (9)--(10) to all net points and using
\[
 \lVert D\rVert_{\mathrm{op}}
 \le2\max_{u\ \text{ in the net}}|u^{\mathsf T}Du|,
 \qquad
 \lVert b\rVert_2
 \le2\max_{u\ \text{ in the net}}|b^{\mathsf T}u|
 \tag{11}
\]
on each support gives, with probability at least \(1-4e^{-x}\),
\[
 d_r:=\sup_{\substack{\lVert u\rVert_2=1\\\lVert u\rVert_0\le r}}
 |u^{\mathsf T}(M-\Sigma)u|
 \le C_3\left(\sqrt{\frac{qL_r}{m}}+\frac{L_r}{m}\right),
 \tag{12}
\]
\[
 b_r:=\sup_{\substack{\lVert u\rVert_2=1\\\lVert u\rVert_0\le r}}
 |\bar\xi^{\mathsf T}u|
 \le C_3\left(\sqrt{\frac{qL_r}{m}}+\frac{L_r}{m}\right),
 \qquad L_r=r\log(en/r)+x.
 \tag{13}
\]

We shall also use the following deterministic consequence of (12)--(13).  If \(D\) is symmetric, decreasing rearrangement of the coordinates of \(v\) into blocks of size \(\lfloor r/2\rfloor\), followed by polarization on every pair of blocks, gives
\[
 |v^{\mathsf T}Dv|
 \le C_4d_r\left(\lVert v\rVert_2+\frac{\lVert v\rVert_1}{\sqrt r}\right)^2.
 \tag{14}
\]
The same blocking without polarization gives
\[
 |b^{\mathsf T}v|
 \le C_4b_r\left(\lVert v\rVert_2+\frac{\lVert v\rVert_1}{\sqrt r}\right).
 \tag{15}
\]
Indeed, monotone ordering bounds the sum of the Euclidean norms of all blocks after the first by \(2\lVert v\rVert_1/\sqrt r\); summing the resulting within-block and cross-block products proves (14)--(15).

Take \(r=\min\{n,\lceil C_5(s+1)\rceil\}\), with \(C_5\) a sufficiently large numerical constant, and set \(x=\log(4/\rho)\).  Since \(s\le n/8\),
\[
 L_r\le C_6\left\{(2s+1)\log\!\left(\frac{en}{2s+1}\right)
                       +\log\!\left(\frac2\rho\right)\right\}.
 \tag{16}
\]
The assumed lower bound on \(m\) makes the right sides of (12)--(13), respectively, at most an arbitrarily small universal multiple of \(q\) and of \(\sqrt q\).

It remains to compare with the population.  By the exact-budget covariance lemma,
\[
 \Sigma=p(1-p)\frac{n}{n-1}
 \left(I_n-\frac1n\mathbf 1_n\mathbf 1_n^{\mathsf T}\right).
 \tag{17}
\]
Because \(p(1-p)=q(1-q)\ge q/2\), the sparse consequence of that lemma and \(2s+1<n\) give
\[
 v^{\mathsf T}\Sigma v\ge c_1q\lVert v\rVert_2^2
 \quad\text{if }\lVert v\rVert_0\le2s+1.
 \tag{18}
\]
For a cone vector, write \(k=|S|\), \(N=n-k\), \(u=\lVert v_S\rVert_1\), and \(w=\lVert v_{S^c}\rVert_1\).  If \(w\le a_0u\) and \(N>a_0k\), Cauchy--Schwarz separately on the two blocks and maximization over \(0\le w/u\le a_0\) give
\[
 \frac{(\mathbf 1_n^{\mathsf T}v)^2}{n\lVert v\rVert_2^2}
 \le \frac{(1+a_0)^2kN}{n(N+a_0^2k)}.
 \tag{19}
\]
For \(a_0=1\), \(k\le s+1\) and \(s\le n/8\) make (19) at most \(3/4\), including \(n=8,s=1\).  Hence (17) gives
\[
 v^{\mathsf T}\Sigma v\ge(q/8)\lVert v\rVert_2^2.
 \tag{20}
\]
For \(a_0=3\), the additional condition \(n\ge9\) gives \(k/n\le2/9\) and \(N>3k\); the right side of (19) is then at most \(224/225\).  Thus
\[
 v^{\mathsf T}\Sigma v\ge(q/450)\lVert v\rVert_2^2.
 \tag{21}
\]
The respective cone conditions also give \(\lVert v\rVert_1\le(1+a_0)\sqrt{s+1}\lVert v\rVert_2\).  Equations (12)--(16), with their universal constants chosen small relative to (18), (20), and (21), therefore transfer all three population lower bounds to the sample matrix because
\[
 \frac1mX^{\mathsf T}X=M-\bar\xi\bar\xi^{\mathsf T}.
 \tag{22}
\]
Finally, \(e_j^{\mathsf T}\Sigma e_j=p(1-p)\le q\).  Equation (12) for one-sparse vectors and (22) imply
\[
 m^{-1}\lVert Xe_j\rVert_2^2\le2q
\]
simultaneously for all \(j\), after a final universal enlargement of \(C_0\).  This proves every asserted conclusion.